% GENERATED FILE. Edit canonical lessons, then run npm run generate:books.
% canonical-source-hash: fnv1a64:ac3e4055055bf220
% canonical-lessons: TE-C192-vatti, TE-C192-gollem, TE-C192-kunce, TE-C192-haram, TE-C192-jada, TE-R192-first-pass-words-from-chapters-184-187, TE-R192-second-pass-words-from-chapters-188-192

\chapter{Wick, Latch, Paintbrush, Necklace, Braid}
\label{ch:te-wick-latch-paintbrush-necklace-braid}
\hlchaptermodality{192}

\begin{chapteropening}
\textbf{By the end of this chapter:} \emph{I can say a wick, a latch, a paintbrush, a necklace and a braid.}

Complete the last lesson of chapter 192: I can say a wick, a latch, a paintbrush, a necklace and a braid.
\end{chapteropening}

\section[\texorpdfstring{\te{వత్తి} (vatti)}{వత్తి (vatti)}]{\te{వత్తి} (vatti) — a wick}

\label{lesson:TE-C192-vatti}



\begin{quote}
Before the new one: say the Telugu for a net, then the Telugu for a weighing scale.
\end{quote}

\subsection*{You'll want to know: \te{వత్తి}}
\textbf{\te{వత్తి}} — \emph{vatti} — \textquotedblleft{}a wick\textquotedblright{}.

\begin{grammarlens}[title={What you've built}]
One more word for this chapter.
\end{grammarlens}

\subsection*{Your turn}
\begin{itemize}
\raggedright
  \item \emph{Say it:} \emph{vatti}
  \item \emph{Say it:} \emph{vatti}, once more
  \item \emph{Say it:} say \emph{vatti}
  \item \emph{Recall:} say the Telugu for a net, then the Telugu for a weighing scale, then say \emph{vatti} again
\end{itemize}

\begin{tcolorbox}[breakable,colback=teal!4,colframe=teal!35!black,title={Before you move on}]
What does \te{వత్తి} mean? (\textquotedblleft{}A wick\textquotedblright{}.) Say it once more.
\end{tcolorbox}
\section[\texorpdfstring{\te{గొళ్ళెం} (goḷḷeṁ)}{గొళ్ళెం (goḷḷeṁ)}]{\te{గొళ్ళెం} (goḷḷeṁ) — a latch, a bolt}

\label{lesson:TE-C192-gollem}



\begin{quote}
Before the new one: say the Telugu for a weighing scale, then the Telugu for a wick.
\end{quote}

\subsection*{You'll want to know: \te{గొళ్ళెం}}
\textbf{\te{గొళ్ళెం}} — \emph{goḷḷeṁ} — \textquotedblleft{}a latch, a bolt\textquotedblright{}.

\begin{grammarlens}[title={What you've built}]
One more word for this chapter.
\end{grammarlens}

\subsection*{Your turn}
\begin{itemize}
\raggedright
  \item \emph{Say it:} \emph{goḷḷeṁ}
  \item \emph{Say it:} \emph{goḷḷeṁ}, once more
  \item \emph{Say it:} say \emph{goḷḷeṁ}
  \item \emph{Recall:} say the Telugu for a weighing scale, then the Telugu for a wick, then say \emph{goḷḷeṁ} again
\end{itemize}

\begin{tcolorbox}[breakable,colback=teal!4,colframe=teal!35!black,title={Before you move on}]
What does \te{గొళ్ళెం} mean? (\textquotedblleft{}A latch, a bolt\textquotedblright{}.) Say it once more.
\end{tcolorbox}
\section[\texorpdfstring{\te{కుంచె} (kuñce)}{కుంచె (kuñce)}]{\te{కుంచె} (kuñce) — a paintbrush}

\label{lesson:TE-C192-kunce}



\begin{quote}
Before the new one: say the Telugu for a wick, then the Telugu for a latch.
\end{quote}

\subsection*{You'll want to know: \te{కుంచె}}
\textbf{\te{కుంచె}} — \emph{kuñce} — \textquotedblleft{}a paintbrush\textquotedblright{}.

\begin{grammarlens}[title={What you've built}]
One more word for this chapter.
\end{grammarlens}

\subsection*{Your turn}
\begin{itemize}
\raggedright
  \item \emph{Say it:} \emph{kuñce}
  \item \emph{Say it:} \emph{kuñce}, once more
  \item \emph{Say it:} say \emph{kuñce}
  \item \emph{Recall:} say the Telugu for a wick, then the Telugu for a latch, then say \emph{kuñce} again
\end{itemize}

\begin{tcolorbox}[breakable,colback=teal!4,colframe=teal!35!black,title={Before you move on}]
What does \te{కుంచె} mean? (\textquotedblleft{}A paintbrush\textquotedblright{}.) Say it once more.
\end{tcolorbox}
\section[\texorpdfstring{\te{హారం} (hāraṁ)}{హారం (hāraṁ)}]{\te{హారం} (hāraṁ) — a necklace}

\label{lesson:TE-C192-haram}



\begin{quote}
Before the new one: say the Telugu for a latch, then the Telugu for a paintbrush.
\end{quote}

\subsection*{You'll want to know: \te{హారం}}
\textbf{\te{హారం}} — \emph{hāraṁ} — \textquotedblleft{}a necklace\textquotedblright{}.

\begin{grammarlens}[title={What you've built}]
One more word for this chapter.
\end{grammarlens}

\subsection*{Your turn}
\begin{itemize}
\raggedright
  \item \emph{Say it:} \emph{hāraṁ}
  \item \emph{Say it:} \emph{hāraṁ}, once more
  \item \emph{Say it:} say \emph{hāraṁ}
  \item \emph{Recall:} say the Telugu for a latch, then the Telugu for a paintbrush, then say \emph{hāraṁ} again
\end{itemize}

\begin{tcolorbox}[breakable,colback=teal!4,colframe=teal!35!black,title={Before you move on}]
What does \te{హారం} mean? (\textquotedblleft{}A necklace\textquotedblright{}.) Say it once more.
\end{tcolorbox}
\section[\texorpdfstring{\te{జడ} (jaḍa)}{జడ (jaḍa)}]{\te{జడ} (jaḍa) — a braid, a plait (of hair)}

\label{lesson:TE-C192-jada}



\begin{quote}
Before the new one: say the Telugu for a paintbrush, then the Telugu for a necklace.
\end{quote}

\subsection*{You'll want to know: \te{జడ}}
\textbf{\te{జడ}} — \emph{jaḍa} — \textquotedblleft{}a braid, a plait (of hair)\textquotedblright{}.

\begin{grammarlens}[title={What you've built}]
One more word for this chapter.
\end{grammarlens}

\subsection*{Your turn}
\begin{itemize}
\raggedright
  \item \emph{Say it:} \emph{jaḍa}
  \item \emph{Say it:} \emph{jaḍa}, once more
  \item \emph{Say it:} say \emph{jaḍa}
  \item \emph{Recall:} say the Telugu for a paintbrush, then the Telugu for a necklace, then say \emph{jaḍa} again
\end{itemize}

\begin{tcolorbox}[breakable,colback=teal!4,colframe=teal!35!black,title={Before you move on}]
What does \te{జడ} mean? (\textquotedblleft{}A braid, a plait (of hair)\textquotedblright{}.) Say it once more.
\end{tcolorbox}
\section[(dialogue)]{Words from chapters 184-187 — a first pass}

\label{lesson:TE-R192-first-pass-words-from-chapters-184-187}



\begin{quote}
Say the words for a necklace and a braid.
\end{quote}

\subsection*{Your turn}
Say these aloud:

\begin{itemize}
\raggedright
  \item butter, a chutney, a sweet milk pudding, a cucumber, okra
  \item an envelope, a postage stamp, a parcel, the post, a pencil
  \item a grandson, a granddaughter, parents, twins, a baby
  \item a bee, a lizard, an insect, a sparrow, a koel
\end{itemize}

\begin{tcolorbox}[breakable,colback=teal!4,colframe=teal!35!black,title={Before you move on}]
Butter? (\textbf{\te{వెన్న}}, \emph{venna}.) A chutney? (\textbf{\te{పచ్చడి}}, \emph{paccaḍi}.) A sweet milk pudding? (\textbf{\te{పాయసం}}, \emph{pāyasaṁ}.) A cucumber? (\textbf{\te{దోసకాయ}}, \emph{dōsakāya}.) Okra? (\textbf{\te{బెండకాయ}}, \emph{beṇḍakāya}.) An envelope? (\textbf{\te{కవరు}}, \emph{kavaru}.) A postage stamp? (\textbf{\te{స్టాంపు}}, \emph{sṭāmpu}.) A parcel? (\textbf{\te{పార్సెలు}}, \emph{pārselu}.) The post? (\textbf{\te{తపాలా}}, \emph{tapālā}.) A pencil? (\textbf{\te{పెన్సిలు}}, \emph{pensilu}.) A grandson? (\textbf{\te{మనవడు}}, \emph{manavaḍu}.) A granddaughter? (\textbf{\te{మనవరాలు}}, \emph{manavarālu}.) Parents? (\textbf{\te{తల్లిదండ్రులు}}, \emph{tallidaṇḍrulu}.) Twins? (\textbf{\te{కవలలు}}, \emph{kavalalu}.) A baby? (\textbf{\te{పాప}}, \emph{pāpa}.) A bee? (\textbf{\te{తేనెటీగ}}, \emph{tēneṭīga}.) A lizard? (\textbf{\te{బల్లి}}, \emph{balli}.) An insect? (\textbf{\te{పురుగు}}, \emph{purugu}.) A sparrow? (\textbf{\te{పిచ్చుక}}, \emph{piccuka}.) A koel? (\textbf{\te{కోకిల}}, \emph{kōkila}.)
\end{tcolorbox}
\section[(dialogue)]{Words from chapters 188-192 — a second pass}

\label{lesson:TE-R192-second-pass-words-from-chapters-188-192}



\begin{quote}
Say the words for prasadam and rice-flour sweets.
\end{quote}

\subsection*{Your turn}
Say these aloud:

\begin{itemize}
\raggedright
  \item prasadam, rice-flour sweets, a festoon of mango leaves hung over a doorway, a bindi, henna
  \item a competition, a winner, a medal, a hobby, a novel
  \item electricity, a machine, a sofa, a swing, a carpet
  \item rubber, a light bulb, firewood, a net, a weighing scale
  \item a wick, a latch, a paintbrush, a necklace, a braid
\end{itemize}

\begin{tcolorbox}[breakable,colback=teal!4,colframe=teal!35!black,title={Before you move on}]
Prasadam? (\textbf{\te{ప్రసాదం}}, \emph{prasādaṁ}.) Rice-flour sweets? (\textbf{\te{అరిసెలు}}, \emph{ariselu}.) A festoon of mango leaves hung over a doorway? (\textbf{\te{తోరణం}}, \emph{tōraṇaṁ}.) A bindi? (\textbf{\te{బొట్టు}}, \emph{boṭṭu}.) Henna? (\textbf{\te{గోరింటాకు}}, \emph{gōriṇṭāku}.) A competition? (\textbf{\te{పోటీ}}, \emph{pōṭī}.) A winner? (\textbf{\te{విజేత}}, \emph{vijēta}.) A medal? (\textbf{\te{పతకం}}, \emph{patakaṁ}.) A hobby? (\textbf{\te{అభిరుచి}}, \emph{abhiruci}.) A novel? (\textbf{\te{నవల}}, \emph{navala}.) Electricity? (\textbf{\te{కరెంటు}}, \emph{kareṇṭu}.) A machine? (\textbf{\te{యంత్రం}}, \emph{yantraṁ}.) A sofa? (\textbf{\te{సోఫా}}, \emph{sōphā}.) A swing? (\textbf{\te{ఉయ్యాల}}, \emph{uyyāla}.) A carpet? (\textbf{\te{తివాచీ}}, \emph{tivācī}.) Rubber? (\textbf{\te{రబ్బరు}}, \emph{rabbaru}.) A light bulb? (\textbf{\te{బల్బు}}, \emph{balbu}.) Firewood? (\textbf{\te{కట్టెలు}}, \emph{kaṭṭelu}.) A net? (\textbf{\te{వల}}, \emph{vala}.) A weighing scale? (\textbf{\te{త్రాసు}}, \emph{trāsu}.) A wick? (\textbf{\te{వత్తి}}, \emph{vatti}.) A latch? (\textbf{\te{గొళ్ళెం}}, \emph{goḷḷeṁ}.) A paintbrush? (\textbf{\te{కుంచె}}, \emph{kuñce}.) A necklace? (\textbf{\te{హారం}}, \emph{hāraṁ}.) A braid? (\textbf{\te{జడ}}, \emph{jaḍa}.)
\end{tcolorbox}
